Ce rapport présente la conception et l'implémentation complète d'un interpréteur
pour \ilang, un mini-langage de programmation impératif conçu dans le cadre du
projet final du cours d'Interprétation \& Compilation. Le langage supporte les
variables, les expressions arithmétiques avec priorité des opérateurs, les
structures de contrôle (conditions, boucles), les fonctions utilisateur
récursives et les entrées-sorties. Sa réalisation réinvestit l'ensemble de la
chaîne de compilation étudiée en cours : analyse lexicale par automate fini
(Flex), analyse syntaxique par grammaire LALR(1) (Bison), construction d'un arbre
syntaxique abstrait (AST), optimisations sur cet arbre (pliage de constantes,
élimination du code mort), puis évaluation récursive.

Au-delà du cahier des charges, le projet a été prolongé par trois extensions
(nombres flottants, instructions \code{arrêter}/\code{continuer}, localisation
des erreurs à la ligne) et par une exploration hors-périmètre particulièrement
formatrice : un \emph{second back-end} qui, à partir du même AST, génère du code
assembleur i386 32~bits compilable avec \code{gcc -m32}. Cette dualité
interpréteur/compilateur illustre concrètement le principe d'un front-end commun
suivi de back-ends interchangeables.

L'implémentation, écrite en C, compile sans aucun avertissement ni conflit
grammatical et ne présente aucune fuite mémoire (vérifié avec Valgrind). Une
campagne de tests systématique, couvrant les programmes de démonstration, la
gestion des erreurs, les cas limites et les mesures de performance, valide le bon
fonctionnement et
met en évidence les forces et les limites de l'approche, notamment la profondeur
de récursion bornée par la pile de l'hôte et l'écart de performance considérable
entre l'interprétation et le code natif généré.

\vspace{0.5cm}
\noindent\textbf{Mots-clés :} interpréteur, compilation, Flex, Bison, LALR(1),
arbre syntaxique abstrait, table des symboles, pliage de constantes, génération
de code, assembleur i386.
